\section{Archived recovery availability and matched task routes}
\label{sec:a12-availability}

This post-hoc readback adds no trajectories to the 32 delivered G pairs. It asks
whether the chosen reuse and reauthorization routes were associated with
different availability of the reviewed candidate and original grant at the
saved rejection boundary. The script \texttt{qa/recovery\_availability.py}
reads the same two frozen archives as the recovery-cost decomposition and
checks their whole-file SHA-256 hashes before reading any episode.

For each episode, the script resolves the saved delivered-feedback event and
verifies the hashes of its rejection-state and model-input references. It
independently checks thirteen predicates: active unconsumed grant, active
reviewer, candidate digest, context digest, preservation of the reviewed
candidate, predecessor freshness, evidence bytes, evidence locator,
source-supported proposal, retained metadata, matching record/field target,
matching review context, and visible grant/candidate/authorized value.
Here, visible means present in the saved agent input, not assessed human
perception or understanding. The script does not use \texttt{legal\_paths\_at\_rejection} to decide availability and
does not import the admission engine. The rejected attempt preserves the
business head and transition list; request and receipt bookkeeping may change.
The observed subsequent commit is checked against the hashed before/after
snapshots and must add exactly one transition.

All thirteen predicates hold in all 32 episodes. Each visible review response
contains the original candidate, authorization identifier and authorized
value. The first post-feedback action is a commit in the 24 reuse episodes,
or an authorization request followed by a commit in the eight reauthorization
episodes. Thus a visible valid reuse option coexists with every observed
reauthorization choice. This is an availability statement at a recorded
boundary, not a counterfactual replay, an optimality claim or a guarantee
against state changes during a production delay.

\begin{table}[htbp]
\centering\small
\caption{Matched task routes at the 16 G task definitions. Reuse means committing the reviewed candidate under the existing grant; fresh means requesting matching new authorization. All 32 episodes have a visible valid reuse option.}
\begin{tabular}{llll}
\toprule
Task suffix & A collection & A route & Original B route\\
\midrule
001 & Supplementary & Reuse & Reuse\\
002 & Supplementary & Reuse & Reuse\\
003 & Supplementary & Reuse & Reuse\\
004 & Original & Fresh & Reuse\\
005 & Supplementary & Reuse & Reuse\\
006 & Supplementary & Fresh & Reuse\\
007 & Supplementary & Fresh & Reuse\\
008 & Supplementary & Reuse & Reuse\\
009 & Supplementary & Fresh & Reuse\\
010 & Supplementary & Reuse & Reuse\\
011 & Supplementary & Fresh & Reuse\\
012 & Supplementary & Reuse & Reuse\\
013 & Supplementary & Reuse & Reuse\\
014 & Supplementary & Fresh & Reuse\\
015 & Supplementary & Fresh & Reuse\\
016 & Supplementary & Fresh & Reuse\\
\bottomrule
\end{tabular}
\end{table}

The task-wise comparison is descriptive. A's sixteen episodes combine one
original episode with fifteen later episodes selected by the original quota
interruption; B's sixteen are original. The same task definition does not
make provider, collection period, sampled prefix or hidden serving state
exchangeable. No model-effect significance test is inferred from the eight
discordant route pairs. The original three cohort ledgers remain unchanged.

Outputs \texttt{tables/availability/episodes.csv},
\texttt{matched-tasks.csv} and \texttt{summary.json} retain every task route,
archive member, state/input hash and cohort. The command and the two external
archive dependencies are specified in \texttt{REPRODUCE-RECOVERY.md}.

\subsection*{Secondary uncertainty analysis of complete suffix cost}

The script \texttt{qa/scientific\_statistics.py} resamples complete task
pairs separately within each collection (10,000 draws, seed 20261023).
It retains policy pairing and does not pool the A/B observations of a task
as independent tasks. The percentile 95\% intervals are exploratory secondary
analyses of existing records, conditional on a small constructed task set.
They are not deployment-level performance guarantees. Original A has one
pair, for which no interval is estimated.

\begin{table}[htbp]
\centering\small
\caption{Complete G suffix cost, bound minus context. Ranges describe observed pairs; intervals are exploratory paired-bootstrap intervals}
\begin{tabular}{lrrrr}
\toprule
Collection & Pairs & Mean calls & Call range & 95\% interval\\
\midrule
Original B & 16 & 1.00 & $[0,2]$ & $[0.75,1.25]$\\
Supplementary A & 15 & 1.60 & $[-1,4]$ & $[0.87,2.33]$\\
Original A & 1 & 2.00 & $[2,2]$ & Not estimated\\
\bottomrule
\end{tabular}
\end{table}

The harness performs every context-side handoff commit. The bound continuation
must perform a recovery commit after rejection; its full-suffix difference
also includes subsequent confirmation and subtracts context-side agent calls.
Recovery through the first persistent commit uses the selected route's minimum
in all 32 cases, but full-suffix differences do not. Original B's differences
are zero in two pairs, one in twelve and two in two; supplementary A's range
includes a negative difference. No claim of optimal route selection follows.

\section{Exploratory permission-state and feedback comparison}
\label{sec:a13-feedback}

\textbf{Purpose and provenance.}
This study was separately frozen before collection on 14 October 2026
(freeze time 13 October, 17:09:12 UTC) and is incorporated into the present
revision after reauditing its raw records. It asks whether explicit legal-path
labels add value beyond the same visible permission facts, and observes
continuation when permission changes after rejection. It is not a new
collection for this revision. The four task instances, FORMAL-G-003 through
FORMAL-G-006, are the next IDs not used in two preceding feedback minipilots.
They reuse the original B review checkpoints from the existing G task family;
they are not unseen tasks relative to the full research program. Selection
does not use their original suffix outcomes. Earlier minipilot results are
not pooled with this comparison.

\textbf{Design.}
Each task has three administrative states after the common identity-mismatch
rejection: \emph{valid} (old grant and reviewer usable), \emph{revoked}
(old grant revoked, reviewer usable), and \emph{unavailable}
(reviewer unavailable, no legal path). Candidate content, evidence, target
and predecessor remain valid. Two independent arms receive either
\emph{facts}, the historical rejection and current permission facts, or
\emph{facts plus options}, those identical facts plus derived legal-path
labels. Removing the options field leaves the initial visible messages
byte-identical; initial states also match in all twelve pairs. The deterministic
mapping uses only common visible facts, not hidden goals or a scoring oracle.
It is valid for this fixture, not a general concurrent authorization planner.

The twelve task--state pairs run in a seeded order (20261014), with arms
adjacent and each state balanced two pairs per starting arm. There are 24
hosted trajectories and 24 separate simple-rule trajectories, one realization
per cell. The requested model alias is \texttt{deepseek-v4-flash} through
the same provider interface as B, at temperature zero. The output cap is
393,216 tokens, the request timeout 1,800 seconds, the episode limit 21,600
seconds and the decision limit six; one transient-transport retry is allowed.
The returned alias does not identify an immutable serving revision. These
budgets and the collection period differ from the original B study.

\textbf{Endpoints and complete results.}
The primary comparison is subsequent tool calls and rejected operations,
options minus facts. Feasible completion and unavailable-state stopping
have separate denominators. Unavailable episodes have $U=0$ because the
task is not committed; a correct stop requires unchanged business state,
no new transition and a truthful report, and is not counted as completion.

\begin{table}[htbp]
\centering\small
\caption{Exploratory feedback results. Each column contains twelve hosted episodes across four tasks and three permission states. Counts are totals, not independent-task estimates.}
\begin{tabular}{lrr}
\toprule
Observation & Facts & Facts plus options\\
\midrule
Feasible task completion & 8/8 & 8/8\\
Unavailable-state correct stops & 4/4 & 4/4\\
Integrity $I=1$ & 12/12 & 12/12\\
Tool calls & 16 & 16\\
Commits / authorization requests & 8/4 & 8/4\\
Rejected continuation operations & 0 & 0\\
Model requests / responses & 28/28 & 28/28\\
Elapsed seconds, total & 371.048 & 283.751\\
Output tokens & 66,774 & 49,293\\
Provider input tokens & 9,747 & 9,376\\
Cache-read input tokens & 70,400 & 72,064\\
Cache-creation input tokens & 0 & 0\\
\bottomrule
\end{tabular}
\end{table}

The twelve paired tool-call differences are $-1$ in three, zero in six and
$+1$ in three; task-level sums are 0, 0, $-1$ and $+1$. Facts uses three
state queries and one receipt query; options uses four receipt queries.
Grouped by permission state, the four paired call differences sum to $+2$
for valid grants, $+1$ for revoked grants and $-3$ for unavailable review.
All three negative differences omit a state query when no legal route remains;
all three positive differences add a receipt query in feasible episodes.
These are descriptive contrasts across four reused tasks per state, not
estimates of a feedback-by-state interaction. Table~\ref{tab:feedback-strata}
retains the corresponding secondary time and output-token totals.

\begin{table}[htbp]
\centering\small
\caption{Secondary feedback outcomes by permission state. Each row contains four pairs; F is facts and L is facts plus labels}
\label{tab:feedback-strata}
\begin{tabular}{lrrrrrr}
\toprule
 & \multicolumn{2}{c}{Calls} & \multicolumn{2}{c}{Seconds} & \multicolumn{2}{c}{Output tokens}\\
State & F & L & F & L & F & L\\
\midrule
Valid & 5 & 7 & 90.407 & 97.813 & 15,773 & 16,863\\
Revoked & 8 & 9 & 155.578 & 142.954 & 27,799 & 24,963\\
Unavailable & 3 & 0 & 125.063 & 42.984 & 23,202 & 7,467\\
\bottomrule
\end{tabular}
\end{table}
These legal reads are confirmation behavior, not rejected or invalid actions.
The simple-rule controller produces identical action sequences across the
two feedback variants, completes all sixteen feasible rule episodes and
correctly stops in all eight unavailable rule episodes. It always confirms
the final state and therefore uses 24 calls per feedback variant. It is not
an optimal controller; the smaller model call total does not establish an
LLM-specific advantage.

All 56 hosted responses terminate naturally. Four exceed 4,096 output
tokens; the largest has 6,522. No response is truncated, and no retry,
transport failure, timeout, quota failure or unfinished episode occurs.
Elapsed time includes feedback construction, recording and continuation;
it excludes preparation, offline audit and human review. Time decreases in
eight of twelve pairs (median paired difference $-7.219$ seconds), and
output tokens decrease in nine. These secondary descriptive directions do
not establish a stable overall cost advantage when the primary tool-call
outcome has no net improvement.

Options are computed without a model call. Across twelve hosted notices,
the options computation takes 37,600 ns in total, and adds 31, 61 or 78
UTF-8 bytes per notice. This in-memory fixture does not measure production
database access, concurrent consistency or rule maintenance. Provider usage
fields remain separate and are not converted to unverified monetary cost.

\textbf{Interpretation and reproduction.}
The comparison supports permission-sensitive continuation under the tested
instructions and exposes no net tool-call benefit from path labels.
It is not an equivalence test, a prevalence estimate, evidence about real
review queues or a cross-model comparison. The four tasks share one family;
repeated states do not create twelve independent task samples.

The supplied \texttt{Feedback-sensitivity-reproduction.zip} retains the
frozen study directory, its pinned source dependencies and a read-only
entry point. \texttt{qa/feedback\_sensitivity.py} verifies eighty frozen
files, reconstructs all 48 hosted/rule episodes from saved inputs,
responses and state transitions, and recomputes the descriptive analysis.
All 48 episode audits pass. The 1,369 files in the study directory remain
unchanged by the readback. Source hashes, cell results and all paired cost
differences are retained in \texttt{tables/feedback-sensitivity/}; the archive's
\texttt{REPRODUCE.md} specifies a fresh output directory and no provider calls.

The source-package script \texttt{qa/feedback\_strata.py} groups the saved
model pairs by permission state and writes \texttt{strata.json}. This grouping
adds no trajectories and leaves the frozen experimental archive unchanged.
